% Ce chapitre repose sur :
%   DOC: s01, s03, s05, s12, s20, s27
%   DOC: s28, s29, s30
%   OBS: REL-DEM.P02, REL-DEM.P03
%   HYP: (aucun)
%   CHI: capacite-annoncee, capacite-fonctionnelle, missions-service, profils-applicatifs
%   SER: (aucune)
%   SECTIONS: 2

\chapter{L'organisme d'accueil}\label{chap:organisme}

L'hôpital provincial Sidi Saïd de Meknès a accueilli le stage, dans son service d'accueil et
d'admission. Quelle place l'hôpital occupe-t-il dans l'offre de soins publique, et quelles missions
le règlement confie-t-il à ce service ?

\section{L'hôpital Sidi Saïd}

L'hôpital appartient à la catégorie des hôpitaux provinciaux ou préfectoraux, la plus nombreuse du
réseau hospitalier public \cite[tableau~32, p.~46]{s30}. À ce niveau, la fonction d'urgence est
assurée sous sa forme de base \cite{s12}. L'hôpital se situe dans la région Fès-Meknès \cite{s05}.
Il dépend du centre hospitalier préfectoral Mohammed~V, dont il est l'une des cinq structures
\cite{s20, s29}. Ce rattachement fixe le niveau auquel les grandeurs sont publiées. Certaines
statistiques nationales descendent jusqu'à l'hôpital nommé, d'autres s'arrêtent au centre
hospitalier ou à la préfecture, et le jeu de données en dépend (chapitre~\ref{chap:jeu-de-donnees}).

L'hôpital a été rénové en 2019, avec la création d'un pôle Mère-Enfant \cite{s29}. La
figure~\ref{fig:organisation-etablissement} montre ses départements et ses services, dont le service
d'accueil et d'admission.

\begin{figure}[htbp]
\centering
\begin{tabular}{@{}c@{}}
\fbox{\begin{minipage}{0.90\linewidth}
\centering\small
\fbox{\begin{minipage}{0.72\linewidth}\centering Centre hospitalier préfectoral Mohammed~V \\ \footnotesize cinq structures\end{minipage}}\\[4pt]
$\downarrow$\\[4pt]
\fbox{\begin{minipage}{0.72\linewidth}\centering \textbf{Hôpital provincial Sidi Saïd}\end{minipage}}\\[6pt]
\fbox{\begin{minipage}{0.30\linewidth}\centering Département \\ de médecine\end{minipage}}%
\hfill
\fbox{\begin{minipage}{0.30\linewidth}\centering Département \\ mère-enfant\end{minipage}}%
\hfill
\fbox{\begin{minipage}{0.30\linewidth}\centering Services \\ médico-techniques\end{minipage}}\\[6pt]
\fbox{\begin{minipage}{0.46\linewidth}\centering Consultations externes\end{minipage}}%
\hfill
\fbox{\begin{minipage}{0.46\linewidth}\centering \textbf{Service d'accueil} \\ \textbf{et d'admission} \\ \footnotesize lieu du travail\end{minipage}}
\end{minipage}}
\end{tabular}
\caption[Organisation de l'établissement]{Organisation de l'établissement, telle que la source la décrit \cite{s29}}\label{fig:organisation-etablissement}
\end{figure}

Deux capacités litières sont attribuées à l'hôpital, et elles ne comptent pas les mêmes lits. Une dépêche
d'agence de février~2020 annonce \chiffre{capacite-annoncee} lits \cite{s29}. Le recueil
statistique du ministère relève, pour l'exercice~2024, une capacité litière \emph{fonctionnelle} de
\chiffre{capacite-fonctionnelle} lits \cite{s30}, celle des lits effectivement armés. Plusieurs
titres de presse ont repris la dépêche mot pour mot \cite{s01}, et la valeur annoncée ne repose que
sur elle \cite{s29}.

L'hôpital n'exerce aucune activité chirurgicale. Le recueil statistique ne lui attribue aucune
intervention chirurgicale sur les deux exercices publiés \cite{s30}. La Cour des comptes décrit sa
gynéco-obstétrique comme une maison d'accouchement \cite{s20}. Le
chapitre~\ref{chap:jeu-de-donnees} démontre ce constat et en tire une décision de modélisation.

\section{Le Service d'Accueil et d'Admission}

Le règlement intérieur des hôpitaux régit le fonctionnement du service. Il est fixé par l'arrêté
du ministre de la santé n\textsuperscript{o}~456-11 du 6~juillet 2010 \cite{s03}, publié au
Bulletin Officiel n\textsuperscript{o}~5926 du 17~mars 2011 \cite{s28}. Son article~27 organise
l'hôpital en départements et en services selon sa tranche de capacité litière, sans nommer ces
services \cite{s27}. Il ne décrit pas l'organisation interne du service d'accueil et d'admission.
L'article~35 fixe les neuf missions du service \cite{s27}, que le tableau~\ref{tab:missions-service}
reproduit dans les termes du texte.

\begin{table}[htbp]
\centering
\caption[Les missions du service d'accueil et d'admission]{Les missions du service d'accueil et d'admission, dans les termes du règlement \cite{s27}}\label{tab:missions-service}
\begin{tabular}{@{}r@{\hspace{2mm}}p{12.4cm}@{}}
\toprule
\multicolumn{2}{@{}l@{}}{Les neuf activités dont le service d'accueil et d'admission est chargé} \\
\midrule
1 & Gérer l'accueil et l'orientation des patients \\
2 & Organiser les admissions et les sorties des malades ainsi que la gestion des rendez-vous \\
3 & Gérer les effectifs des patients et leurs mouvements à l'intérieur de l'hôpital \\
4 & Enregistrer les actes médicaux légaux et assurer la gestion de la morgue hospitalière \\
5 & \textbf{Établir les statistiques et gérer l'information hospitalière} \\
6 & \textbf{Établir la facturation des prestations et services rendus par l'hôpital sur la base de
la classification des maladies, des nomenclatures des actes et des tarifs en vigueur} \\
7 & Organiser la communication interne et externe de l'hôpital \\
8 & Assurer l'assistance sociale aux patients \\
9 & Organiser et gérer les archives médicales et les dossiers des patients de l'établissement \\
\bottomrule
\end{tabular}
\end{table}

Le tableau porte en gras les deux missions qui gouvernent ce travail. La sixième fait du service le
lieu où naît la donnée de facturation. La cinquième confie au service lui-même la production des
statistiques. Un tableau de bord alimenté par ses données sert donc l'une de ses attributions.

\medskip

Ce que le service fait au quotidien se lit aussi dans l'outil qu'il emploie. L'écran de connexion
du système d'information propose cinq profils applicatifs, et chacun ouvre un jeu de modules
distinct\releve{REL-DEM.P02}. Les profils ne sont donc pas cinq vues d'un même écran : ils
découpent l'activité du service en cinq processus.

\textbf{Un profil du logiciel sert une mission que le règlement place ailleurs.} Le profil
\texttt{MSM - RECOUVREMENT}\releve{REL-DEM.P03} porte le suivi des créances et des relances, et le
recouvrement ne figure nulle part à l'article~35 : c'est l'article~9, paragraphe~b, qui charge le
pôle des affaires administratives de veiller au recouvrement des créances de l'établissement
\cite{s27}. Le découpage du logiciel ne recopie donc pas le découpage réglementaire, et un service
peut porter à l'écran une activité que le texte attribue à un autre. Ce travail observe le
logiciel tel qu'il est, et cite le règlement tel qu'il est écrit.

L'observation conduite pour ce travail a porté sur un poste et sur un seul de ces cinq profils,
celui des rendez-vous : les quatre autres exigent des habilitations qu'un stagiaire ne reçoit
pas.

Le chemin d'un patient relie ces processus entre eux (figure~\ref{fig:parcours-patient}).

\begin{figure}[htbp]
\centering
\begin{tikzpicture}[x=1mm, y=1mm]

\node[briqueUML] (arr) at (0,0) {Arrivée du patient};
\node[briqueUML] (rdv) at (31,14) {Prise de rendez-vous\\[1pt]\emph{gestion de RDV}};
\node[briqueUML] (urg) at (31,-14) {Passage aux urgences\\[1pt]\emph{urgence}};
\node[briqueUML] (cons) at (62,14) {Consultation externe};
\node[briqueUML] (tri) at (62,-14) {Tri et orientation};
\node[briqueUML] (adm) at (93,14) {Admission\\en hospitalisation};
\node[briqueUML] (sej) at (124,14) {Séjour\\[1pt]\emph{mouvements internes}};
\node[briqueUML] (sor) at (124,-14) {Sortie};
\node[briqueUML] (fac) at (93,-38) {Facturation\\[1pt]\emph{facturation SAA}};
\node[briqueUML] (cre) at (62,-38) {Créance impayée};
\node[briqueUML] (rel) at (31,-38) {Relance\\[1pt]\emph{recouvrement}};

\draw[fluxUML] (arr.east) -- (rdv.west);
\draw[fluxUML] (arr.east) -- (urg.west);
\draw[fluxUML] (rdv) -- (cons);
\draw[fluxUML] (urg) -- (tri);
\draw[fluxUML] (cons) -- (adm);
\draw[fluxUML] (tri.east) -- (adm.south west);
\draw[fluxUML] (adm) -- (sej);
\draw[fluxUML] (sej) -- (sor);

% Les trois voies qui aboutissent à une facture se rejoignent en un point avant elle.
\coordinate (j) at (93,-26);
\fill[accentgarde] (j) circle (0.6mm);
\draw[fluxUML] (cons.south) -- ++(0,-6) -| (j);
\draw[fluxUML] (tri.south) -- ++(0,-6) -| (j);
\draw[fluxUML] (sor.south) -- ++(0,-6) -| (j);
\draw[fluxUML] (j) -- (fac.north);

\draw[fluxUML] (fac) -- (cre);
\draw[fluxUML] (cre) -- (rel);
\end{tikzpicture}
\caption{Le parcours d'un patient à travers les \chiffre{profils-applicatifs} processus du service}\label{fig:parcours-patient}
\end{figure}

Le parcours va de l'arrivée du patient à la relance d'une créance. La voie programmée commence par
un rendez-vous, la voie non programmée par un passage aux urgences. Une consultation, une
orientation de sortie ou une sortie d'hospitalisation aboutissent toutes trois à une facture. Une
facture non soldée devient une créance, et seule une créance appelle une relance. Les intitulés en
italique nomment le processus, non le profil applicatif.

\section*{Synthèse}

L'hôpital Sidi Saïd est un hôpital provincial sans activité chirurgicale, rattaché au centre
hospitalier préfectoral Mohammed~V. Le règlement confie à son service d'accueil et d'admission
\chiffre{missions-service} missions, dont la facturation et la production des statistiques. Le
logiciel découpe ce travail en \chiffre{profils-applicatifs} profils, selon un autre découpage. Le
chapitre~\ref{chap:systeme} décrit ce logiciel et ce que l'observation au poste en a relevé.
